% Draft response to the referee (OMS submission, DTAP manuscript).
% Build: ./build_response.sh  (compiles the revised manuscript first, so that every
% section, table, equation and page number below is read from its .aux file).
\documentclass[11pt]{article}
\usepackage{fontspec}
\usepackage[a4paper,margin=2.4cm]{geometry}
\usepackage{amsmath,amssymb,xcolor,enumitem,xr}
\usepackage[hidelinks]{hyperref}
\externaldocument[ms-]{newref/new}

% Internal cross-references to the change inventory of DTAP_changes_marked.pdf.
% Remove them before sending:  \renewcommand{\internal}[1]{}
\newcommand{\internal}[1]{{\color{gray}\footnotesize\,[#1]}}
% Notes for the co-authors (Vietnamese); remove before sending.
\newcommand{\confirm}[1]{{\color{red!70!black}\footnotesize\,[CẦN XÁC NHẬN: #1]}}

\newcommand{\R}{\mathbb{R}}
\newcommand{\cvx}{\mathrm{cvx}}
\newcommand{\secref}[1]{Section~\ref{ms-#1}}
\newcommand{\pg}[1]{p.~\pageref{ms-#1}}
\newenvironment{comment}[1]{\par\medskip\noindent\textbf{#1}\par\nopagebreak\begin{quote}\itshape}{\end{quote}}
\newenvironment{response}{\par\noindent\textbf{Response.}\ }{\par}
\newcommand{\changes}{\par\noindent\textbf{Changes.}\ }

\begin{document}

\begin{center}
{\Large\bfseries Response to the Referee}\\[4pt]
\emph{A Dual-Track Adaptive Projection Algorithm for a Nonlinear Split Feasibility Problem\\ with Mixed Convex--Pseudoconvex Constraints}\\[2pt]
Nguyen Van The, Tran Ngoc Thang and Nguyen Hong Son
\end{center}

\noindent Dear Editor, dear Referee,

\medskip
\noindent We thank the referee for the careful reading of our manuscript and for the
constructive comments. We have revised the paper substantially. Below, the referee's
comments are reproduced in italics, followed by our response and the corresponding
changes. Section, equation, table and page numbers refer to the revised manuscript.
A version of the manuscript with all changes marked (additions highlighted, deletions
struck through) is provided separately.

\subsection*{Summary of the main changes}
\begin{enumerate}[leftmargin=1.8em,itemsep=2pt]
\item The Introduction and the Related work section were merged into a single, shorter
  Introduction (\secref{sec:intro}). The manuscript was shortened from 39 to 25 pages.
\item The paper was rewritten throughout in shorter sentences. Redundant motivation,
  research questions, lists of difficulties and several subsections were removed.
\item The format of the references was unified, and four works on self-adaptive and
  decomposition methods for the SFP were added. References that are no longer needed
  after shortening were removed.
\item The appendix was removed. The per-dataset results are now in the main text
  (Table~\ref{ms-tab:acc}), and the baseline solvers are described in
  \secref{sssec:baselines}.
\item In addition to the referee's comments, we corrected and sharpened several
  statements, revised the experimental section and added a new solver diagnostic
  (\secref{ssec:feasibility-audit}). These changes are listed in the final section of this
  letter.
\end{enumerate}

% =====================================================================
\section*{Major issues}

\begin{comment}{Comment M1.}
Section 1 ``Introduction'' and Section 2 ``Related work'' are suggested to be
reorganized. Both sections are suggested to be merged as Section 1. Furthermore, the
Related work is unnecessary to be revisited in detail. After merging the Introduction
and Related work, it is suggested to limit them to 4 pages. [\dots] Please shorten the
paper as much as possible. [\dots] Maybe 25 pages is enough for this paper by removing
redundant content and shortening long sentences.
\end{comment}
\begin{response}
We agree and have followed this suggestion. The Related work section has been removed
and its essential content condensed into the Introduction. The revised Introduction
(\secref{sec:intro}, pp.~\pageref{ms-sec:intro}--\pageref{ms-sec:prelim}) recalls the
$CQ$ method, self-adaptive and decomposition strategies, and the closest nonlinear
antecedents (Li et al., Yen et al., Hu et al.) in a few short paragraphs, and it states
three contributions. The whole paper was shortened from 39 to 25 pages.
\end{response}
\changes
\begin{itemize}[leftmargin=1.5em,itemsep=1pt]
\item Introduction rewritten and shortened; the former Section~2 (Related work,
  including its Table~1) removed.\internal{C3--C19}
\item Preliminaries reduced to the facts used later, without subsections
  (\secref{sec:prelim}).\internal{C20--C28}
\item The former Section~5 (feature-selection framework) condensed into
  \secref{sec:fs} of the experiments.\internal{C61--C80}
\item Removed: the sparse variant DTAP-S and its experiments, the runtime comparison,
  the section ``Further analyses of robustness, fidelity, and ablation'' and the
  ``Summary of findings''.\internal{C70, C107--C110}
  \confirm{lý do bỏ DTAP-S và so sánh thời gian chạy --- chỉ do rút gọn, hay còn lý do khác cần nêu?}
\item Conclusion reduced to one paragraph (\secref{sec:conclusion}).\internal{C111--C115}
\end{itemize}

\begin{comment}{Comment M2.}
In order to make the paper more readable, the writing of the current version of the
paper should be substantially improved. Many sentences are ambiguous. [\dots] Please
describe the problem, algorithms, etc.\ in short sentences as much as possible. [\dots]
\end{comment}
\begin{response}
We have rewritten the manuscript with this request in mind. Long sentences were split,
ambiguous phrases were replaced or removed, and descriptive passages that did not
support the results were deleted. In particular, the abstract, the Introduction, the
Preliminaries, the opening paragraphs of \secref{sec:convergence} and the experimental
section were rewritten. The names in parentheses of lemmas, propositions, remarks and
corollaries were removed to shorten the text.
\end{response}
\changes Abstract;\internal{C1} \secref{sec:intro};\internal{C3--C14}
\secref{sec:prelim};\internal{C20--C28} Sections~\ref{ms-ssec:model-stmt}--\ref{ms-ssec:complexity};\internal{C30--C43}
\secref{sec:experiments};\internal{C71--C110} \secref{sec:conclusion}.\internal{C111}

\begin{comment}{Comment M3.}
The format of the references is inconsistent. Many important references are omitted for
SFP methods. E.g., I remember the authors of [23], [35] also have some work on SFP with
self-adaptive or decomposition methods.
\end{comment}
\begin{response}
Thank you for pointing this out. We now cite the following works of the authors of
[23] (Z.~Li, D.~Han, W.~Zhang) and [35] (Q.~Yang) on self-adaptive and decomposition
methods for the (multiple-sets) SFP:
\begin{itemize}[leftmargin=1.5em,itemsep=1pt]
\item W.~Zhang, D.~Han, Z.~Li, A self-adaptive projection method for solving the
  multiple-sets split feasibility problem, \emph{Inverse Problems} 25 (2009) 115001;
\item W.~Zhang, D.~Han, X.~Yuan, An efficient simultaneous method for the constrained
  multiple-sets split feasibility problem, \emph{Comput.\ Optim.\ Appl.} 52 (2012) 825--843;
\item J.~Zhao, Q.~Yang, Self-adaptive projection methods for the multiple-sets split
  feasibility problem, \emph{Inverse Problems} 27 (2011) 035009;
\item Y.~Yang, Q.~Yang, S.~Zhang, Modified alternating direction methods for the modified
  multiple-sets split feasibility problems, \emph{J.\ Optim.\ Theory Appl.} 163 (2014) 130--147.
\end{itemize}
They are cited in the first paragraph of the Introduction, where self-adaptive steps and
decomposition strategies for the SFP are recalled. The reference list now uses one format:
initials for given names, journal names in italics, volume, issue, pages, year and, where
available, the DOI. References that were only needed for the removed material were deleted.
\confirm{phản biện chê \emph{thiếu} tài liệu, nhưng bản sửa bỏ 18 tài liệu, trong đó có các công trình SFP kinh điển (Xu 2010, Fukushima 1986, Censor et al.\ 2006, Chen et al.\ 2019). Cân nhắc giữ lại vài tài liệu này, hoặc chuẩn bị giải thích nếu phản biện hỏi}
\end{response}
\changes First paragraph of \secref{sec:intro};\internal{C3} reference list.\internal{C121--C123}

\begin{comment}{Comment M4.}
Why do the authors present some numerical results in the Appendix?
\end{comment}
\begin{response}
We agree that the numerical results belong in the main text. The appendix has been
removed. The per-dataset accuracy and AUC values are now reported in
Table~\ref{ms-tab:acc} (\pg{tab:acc}), and the implementation of the two baseline solvers
is described in \secref{sssec:baselines}. The table of calibrated thresholds and the
per-dataset runtime table were removed; the thresholds of every instance are recorded
in the supplementary metadata of the reproducibility package.
\end{response}
\changes \secref{sssec:baselines};\internal{C85--C88} Table~\ref{ms-tab:acc};\internal{C100}
appendix removed.\internal{C118--C120}

% =====================================================================
\section*{Minor issues}

\begin{comment}{Comment m1.}
page 2, line 18, what is the ``reference point''? It seems to be ambiguous.
\end{comment}
\begin{response}
We meant that the self-adaptive rule of L\'opez et al.\ was the starting point of our
step-size design. The sentence has been removed. The rule is now cited, together with
other self-adaptive methods, in the first paragraph of the Introduction.
\end{response}
\changes \secref{sec:intro}, first paragraph.\internal{C3}

\begin{comment}{Comment m2.}
page 2, line 27, ``differing densities'' or ``different densities''
\end{comment}
\begin{response}
``Different densities'' is the intended meaning. The sentence (the radiotherapy example)
has been removed in the shortened Introduction.
\end{response}
\changes Paragraph removed.\internal{C9}

\begin{comment}{Comment m3.}
page 2, line 35, please polish the sentence ``The NSFP is not itself new''
\end{comment}
\begin{response}
The sentence has been replaced by ``Nonlinear SFP and generalized convex feasibility have
several relevant antecedents.''
\end{response}
\changes \secref{sec:intro}, third paragraph.\internal{C5, C9}

\begin{comment}{Comment m4.}
page 2, line 39, ``active index set'' or ``active set''?
\end{comment}
\begin{response}
``Active set'' is meant. The sentence has been removed. In the algorithm we use the term
``approximate active set'' (\secref{ssec:dual-track}).
\end{response}
\changes Paragraph removed.\internal{C9}

\begin{comment}{Comment m5.}
page 2, lines 47-49, ``We ask instead \dots\ make precise'' the sentence is difficult for
me to understand.
\end{comment}
\begin{response}
The sentence has been removed. The relation between our function class and those of
the closest works is now stated in short sentences in the third paragraph of the
Introduction: Li et al.\ use a convex aggregate proximity function, which componentwise
pseudoconvexity does not ensure, and Hu et al.\ treat quasiconvex inequalities, which
include our class.
\end{response}
\changes \secref{sec:intro}.\internal{C5, C9}

\begin{comment}{Comment m6.}
page 3, lines 48-53, the sentences are too long; please use short sentences.
\end{comment}
\begin{response}
This paragraph has been removed. Its main point is now stated in one short sentence:
``We seek feasibility, rather than minimization of every component.''
\end{response}
\changes \secref{sec:intro}, second paragraph.\internal{C4, C11, C12}

\begin{comment}{Comment m7.}
page 3, the authors present 5 contributions of this paper. Please summarize those main
contributions instead of describing too many details.
\end{comment}
\begin{response}
The five detailed contributions have been replaced by three short items C1--C3.
\end{response}
\changes \secref{sec:intro}.\internal{C8, C13}

\begin{comment}{Comment m8.}
page 5, ``Rule (3) has since been extended'', please check if it is grammatically correct.
\end{comment}
\begin{response}
The sentence belonged to the former Related work section, which has been removed
(see M1).
\end{response}
\changes Removed.\internal{C16}

\begin{comment}{Comment m9.}
page 5, line 33, what is the meaning of ``drives their proof is for free''? ``the
nonconvexity sits in the sets''?
\end{comment}
\begin{response}
We meant that in the work of Chen et al.\ the operator is a matrix, so the subgradient
inequality used in their proof holds without additional assumptions, and the
nonconvexity is only in the constraint sets, not in the map. The passage was unclear and
has been removed together with the Related work section.
\end{response}
\changes Removed.\internal{C16}

\begin{comment}{Comment m10.}
page 5, line 35, what is ``plain gradient inequality''?
\end{comment}
\begin{response}
We meant the usual gradient inequality of a convex function,
$f(y)\ge f(x)+\langle\nabla f(x),y-x\rangle$, as opposed to co-coercivity of the
gradient. The passage has been removed.
\end{response}
\changes Removed.\internal{C16}

\begin{comment}{Comment m11.}
page 5, line 44, what is ``the closest predecessor''? Do you mean the seminal work?
\end{comment}
\begin{response}
We meant the most closely related earlier work (Li et al.), not the seminal work. The
passage has been removed. The Introduction now refers to these works as ``relevant
antecedents''.
\end{response}
\changes Removed; \secref{sec:intro}.\internal{C5, C17}

\begin{comment}{Comment m12.}
page 5, line 48, ``step'' or ``step size''?
\end{comment}
\begin{response}
``Step size'' was meant. The sentence has been removed with the Related work section.
\confirm{trong bản mới vẫn còn ``residual-scaled step'', ``safeguarded diminishing step''; nếu muốn có thể thống nhất cách dùng ``step''/``step size'' trong toàn bài}
\end{response}
\changes Removed.\internal{C17}

\begin{comment}{Comment m13.}
page 7, line 31, dtap should be DTAP
\end{comment}
\begin{response}
Corrected throughout. The acronym is now always written ``DTAP''.
\end{response}
\changes All occurrences.\internal{C29, C38, C39}

\begin{comment}{Comment m14.}
page 7, line 39, ``following the general framework of Combettes 10'' should be
``following the general framework of Combettes [10]''
\end{comment}
\begin{response}
Corrected. The original sentence was removed with the Related work section; the
remaining citation now reads ``within the framework of Combettes~[\dots]''
(\secref{ssec:bounded}).
\end{response}
\changes \secref{ssec:bounded}, first sentence.\internal{C19, C51}

\begin{comment}{Comment m15.}
page 9, the subtitles 3.1, 3.2 seem unnecessary. Subsection 3.3 is feature selection,
which is suggested to be merged into the Numerical experiments section.
\end{comment}
\begin{response}
We agree. The Preliminaries (now \secref{sec:prelim}) no longer have subsections. The
preliminary subsection on feature selection has been removed, and the feature-selection
model (formerly Section~5) is now \secref{sec:fs} of the Numerical experiments.
\end{response}
\changes \secref{sec:prelim};\internal{C23, C25, C27} \secref{sec:fs}.\internal{C61--C80}

\begin{comment}{Comment m16.}
page 12, $H:=\R^n$ has been stated in the Preliminaries section; why do we repeat
$H:=\R^n$ once again? Please write the paper as concisely as possible.
\end{comment}
\begin{response}
The repetition has been removed from \secref{ssec:model-stmt}.
\end{response}
\changes \secref{ssec:model-stmt}.\internal{C32}

\begin{comment}{Comment m17.}
page 13, the subtitle of 4.3, what is the meaning of ``standing assumption''?
\end{comment}
\begin{response}
We meant assumptions that are in force throughout the analysis. To avoid ambiguity, the
subsection is now titled ``Assumptions'' (\secref{ssec:assumptions}), and
Assumption~\ref{ms-ass:standing} no longer carries the name ``Standing assumptions''.
\end{response}
\changes \secref{ssec:assumptions}.\internal{C36, C37}

\begin{comment}{Comment m18.}
page 18 line 40, [the series] should be $\sum_{k=0}^{\infty}\|x^{k+1}-x^k\|^2<\infty$.
Please follow the advice of the AMS Math formula style as much as possible when
presenting the formula.
\end{comment}
\begin{response}
Corrected. The series in Proposition~\ref{ms-prop:regularity} is now displayed with
the limits above and below the summation sign. We also checked the other displayed
formulas.
\confirm{đã rà soát các công thức khác theo kiểu AMS chưa? Nếu chưa, bỏ câu cuối}
\end{response}
\changes Proposition~\ref{ms-prop:regularity} (\pg{prop:regularity}).\internal{C55}

\begin{comment}{Comment m19.}
page 20, line 36, ``so no Young splitting is needed''. Please polish the writing as much
as possible.
\end{comment}
\begin{response}
The phrase has been removed from the proof of Corollary~\ref{ms-cor:cvx}.
\end{response}
\changes Proof of Corollary~\ref{ms-cor:cvx}.\internal{C59}

\begin{comment}{Comment m20.}
page 20, line 35, $K\cap\bigcap_{i\in I}\cvx$. I don't think this expression is good.
\end{comment}
\begin{response}
We now write $\Omega=\{x\in K:F_i(x)\le b_i\text{ for all }i\in I^{\cvx}\}$ in
Corollary~\ref{ms-cor:cvx}. In Stage~1 of the proof of Theorem~\ref{ms-thm:strong}, the
intersection is now written as the single set
$\bigcap_{i\in I^{\cvx}}\{x\in K:F_i(x)\le b_i\}$, and the case $I^{\mathrm{ps}}=\emptyset$
concludes directly with $C\subseteq\Omega$.
\end{response}
\changes Corollary~\ref{ms-cor:cvx}; proof of Theorem~\ref{ms-thm:strong}.\internal{C56, C59}

\begin{comment}{General comment.}
I suggest the authors consider the above comments in their revision. Furthermore, it
would be recommended if the authors could make the paper concise.
\end{comment}
\begin{response}
We have followed all comments and shortened the paper from 39 to 25 pages (see M1 and M2).
\end{response}

% =====================================================================
\section*{Additional changes made by the authors}
While revising the paper we also made the following changes, which were not requested by
the referee. We list them for completeness.
\begin{enumerate}[leftmargin=1.8em,itemsep=3pt]
\item \textbf{Convexity of the fractional criterion.} The submitted version stated that the
  ratio $F_3(w)=w^{\top}Qw/\rho^{\top}w$ is pseudoconvex. We now prove that it is in fact
  convex on an open neighbourhood of the simplex when $Q\succeq0$ and $\min_k\rho_k>0$
  (Proposition~\ref{ms-prop:fs-reg}). Hence the feature-selection instance belongs to the
  convex subclass. The assignment of $F_3$ to the pseudoconvex track used in the
  experiments remains admissible, and the all-convex assignment is reported as a
  comparator. This is stated in the abstract, the Introduction, \secref{ssec:reduction} and
  the Conclusion.\internal{C1, C7, C64, C68, C78, C111}
\item \textbf{Domain of the component functions.} The functions are now defined on an open
  neighbourhood $U$ of $K$, and (A2) requires continuous differentiability on $U$. This
  covers fractional functions whose denominators are positive on $K$ but not necessarily on
  the whole space. A short
  paragraph after Assumption~\ref{ms-ass:standing} explains this and why $\Omega$ is
  convex.\internal{C32, C35, C37, C49, C50, C54}
\item \textbf{Scope of the method and of the analysis.} New Remarks~\ref{ms-rem:weight-scope}
  and~\ref{ms-rem:regularity-scope} state that the proportional weight is not claimed to be
  unique or optimal, and that the proof uses local bounds on a compact set although no
  constant is an input of the algorithm. Remark~\ref{ms-rem:rate} now states precisely
  that the $O(N^{-1/3})$ bound concerns only the best convex-track violation, and the
  last stage of the proof of Theorem~\ref{ms-thm:strong} was rewritten more
  explicitly.\internal{C40, C41, C56, C58}
\item \textbf{Experimental protocol.} \secref{sssec:protocol} now describes the protocol
  exactly: min--max scaling is fitted once on the training pool (not within each
  selection fold, as stated in the appendix of the submitted version), the voted masks are
  computed once and reused, the SVM probability calibration uses a fixed seed, and the
  evaluation is exploratory rather than nested. The formula of the Nogueira stability
  index and the statistical procedure (exact signed-rank test, Holm adjustment, Friedman
  test) are given explicitly.\internal{C91--C94}
  \confirm{mô tả cũ ``fit on each fold's training portion'' là sai sót mô tả hay code đã đổi? Response cần nói đúng}
\item \textbf{Recomputed values.} The stability values in Table~\ref{ms-tab:nfeat}, the SVM
  AUC values in Tables~\ref{ms-tab:acc} and~\ref{ms-tab:perf-summary}, and the $p$-values
  in Table~\ref{ms-tab:wilcoxon} were recomputed; Table~\ref{ms-tab:wilcoxon}
  now reports raw and Holm-adjusted $p$-values. Accuracies and all $5$-NN results are
  unchanged. Columns for SG-const and SG-dyn were added.\internal{C99, C100, C102, C104}
  \confirm{nêu nguyên nhân chính xác của việc AUC/ổn định thay đổi (seed hiệu chỉnh SVM? hiệu chỉnh $N_f/(N_f-1)$ của chỉ số Nogueira?)}
\item \textbf{Revised claims.} Statements of superiority (accuracy, speed, robustness) were
  removed or qualified. The revised text states that no pairwise comparison is significant
  after Holm adjustment, that the results show a size--prediction trade-off, and that tuned
  subgradient methods can match or outperform DTAP.\internal{C1, C95--C98, C101, C103, C106, C111}
\item \textbf{New solver diagnostic.} \secref{ssec:feasibility-audit} is new. It reports
  numerical feasibility witnesses, the activity of the third constraint along the
  trajectories, a comparison with subgradient methods whose moduli were tuned (including
  the tuning cost), Movement Libras as an additional instance, and controlled examples
  with a genuinely nonconvex pseudoconvex constraint~\eqref{ms-eq:logball-test}
  (Tables~\ref{ms-tab:activation-tuning}--\ref{ms-tab:dtap-sweep}), together with a
  reproducibility paragraph.\internal{C57, C71, C105}
\item \textbf{Model details.} Implementation details of the feature-selection model are now
  stated: the regularization constants $\delta_Q$ and $\varepsilon_\rho$, the variance
  floor in the fuzzy relation, the tie-breaking rule in the decision matrix, the notation $\Pi$ in the
  projection-neural-network equation (to avoid a clash with the decision matrix $P$), the
  projection of the network output onto the simplex, and the fact that the cumulative selection rule neither enforces a feature
  budget nor preserves feasibility.\internal{C74, C76, C79, C87}
\item \textbf{Smaller corrections.} The cost of the simplex projection is $O(p\log p)$ and
  the storage of one iteration is $O(nrp+p^2)$; a typo $\{x_k\}$ was corrected to
  $\{x^k\}$.\internal{C24, C51, C80}
\item \textbf{Declarations.} The declaration of generative AI use and the data-availability
  statement were updated; a reproducibility package now accompanies the
  paper.\internal{C116, C117}
\end{enumerate}

\bigskip
\noindent We hope that the revised manuscript is now suitable for publication, and we thank the
referee again for the helpful comments.

\medskip
\noindent Sincerely,\\
Nguyen Van The, Tran Ngoc Thang and Nguyen Hong Son

\end{document}
